\chapter*{Abstract}

This thesis presents a novel polynomial expansion technique designed to facilitate detailed and exact analysis of a broad class of nonlinear functions.  The primary motivation for this development is to formulate a polynomial expansion which is suitable for application to nonlinear activation functions often employed in neural networks.  In the study of neural networks, polynomial expansions have been applied to address well-known difficulties in the verifiable, explainable, and secure deployment thereof. Existing approaches span classical Taylor and Chebyshev methods, asymptotics, and many numerical and algorithmic approaches.  While these individually have useful properties such as exact error formulas, monic form, adjustable domain, and robustness to undefined derivatives, there are no approaches that provide a consistent method yielding an expansion with all these properties.

To address this gap, an analytically modified integral transform expansion technique, referred to as AMITE, is developed.  The AMITE technique is a novel expansion via integral transforms modified using a derived criterion for convergence.  The AMITE technique is applied to a broad class of functions arising in nonlinear systems, including the rectified linear unit (ReLU), the hyperbolic secant, the hyperbolic secant squared, and the hyperbolic tangent.  Compared with existing state of the art expansion techniques such as Chebyshev, Taylor Series, Fourier-Legendre, and many modern numerical approximations, AMITE is the first polynomial expansion that can provide six previously mutually exclusive desired expansion properties such as exact formulas for the coefficients and exact expansion errors while remaining robust to undefined derivatives.  The utility of the expansion is demonstrated through a main in-depth case study where its use to address the equivalence testing problem of the Multi-layered Perceptron (MLP) is demonstrated.  In this case study, a blackbox network under test is stimulated and a replicated multivariate polynomial form is efficiently extracted from the response to enable comparison against the original network.  A comparison procedure enabled by the expansions is developed and called a conformance test, i.e., a test to check that a neural network implemented  conforms to an original, designed network.

The general utility of the technique is then further demonstrated through a series of case studies in nonlinear waves and nonlinear circuits, motivated by applications in naval radar systems and signal processing. First, a problem involving waves traveling in shallow water is solved where the mean height of a shallow water wave with an uncertain wavenumber is derived. Further, two problems relating to characterization of biased and unbiased nonlinear amplifiers are solved.  In solving these problems, several families of novel integer sequences are derived and tabulated for the first time.  Throughout, a number of nontrivial relationships among the special functions were derived to enable the results obtained.  These relationships and their proofs are cataloged in the appendix.  We conclude that the AMITE technique presents a new dimension of expansion methods that are suitable for precise analysis of nonlinearities arising in a wide variety of nonlinear circuits and systems, and specifically has noteworthy utility in opening new directions and opportunities for the theoretical analysis and testing of neural networks.
